\section{Note Searching Methods}
\label{research:NoteSearchingMethods}

I set out researching methods for storing files / notes in a way that is easily searchable and accessible.

\subsection*{LATCH Method}
\textbf{Source:} \cite{LATCHInformationArchitecture}

The LATCH system, theorised by Richard Saul Wurman in his book *Information Anxiety*, aims to make finding data as simple as possible by using five sorting methods:

\begin{itemize}
  \item \textbf{Location:} Record where the note was physically made, associating memories with a natural link to physical location.
  \item \textbf{Alphabet:} Traditional methods of searching files by their title name/content. Efficient when information is known — where most note systems stop.
  \item \textbf{Time:} Record when notes were created. Like location, this helps the brain retrieve knowledge via natural links.
  \item \textbf{Category:} Use tags and links to relate notes together. Helps the mind bridge from related, even unknown, information.\\
  \textit{→ This is where I want to integrate my AI tools to automate creating these links.}
  \item \textbf{Hierarchy:} Traditional folder/subpage system. Useful when links are obvious but can hide information and require a complex index.
\end{itemize}


\subsection*{PARA Method}
\textbf{Source:} \cite{PARAMethodSimple}

Primarily used for file/project delivery:

\begin{itemize}
  \item \textbf{Projects:} Notes currently in progress or under active editing.
  \item \textbf{Areas:} Files accessed regularly but not being edited (e.g. finances).
  \item \textbf{Resources:} Reference materials such as articles.
  \item \textbf{Archives:} Completed or paused work.
\end{itemize}


\subsection*{Johnny Decimal System}
\textbf{Source:} \cite{SystemOrganiseYour2025}
\begin{enumerate}
  \item Create a maximum of 10 large “buckets” (0–9).
  \item Divide each bucket into 10 subfolders (0–9).
  \item Place each file into the system using a unique ID number.
\end{enumerate}

Each file is referenced like: \texttt{7.3.18} — the 18th file in the 3rd subfolder of the 7th bucket.


\subsection*{eACCESS Framework}
\textbf{Source:} \cite{Lession5Why}

\begin{itemize}
  \item \textbf{Encounters:} Inbox of unsorted items.
  \item \textbf{Atlas:} High-level overviews and indexes (maps of content).
  \item \textbf{Calendar:} Logging, tracking, and reviewing time-based notes.
  \item \textbf{Cards:} Connect and develop ideas.
  \item \textbf{Extras:} Support materials like templates and internal references.
  \item \textbf{Sources:} External knowledge bases.
  \item \textbf{Spaces:} Domains of life (e.g. work, home).
\end{itemize}


\subsection*{Flat Structure}
\textbf{Source:} \cite{staffFlatFileSystem}
All files/folders are stored at the same level — no subfolders.

\begin{itemize}
  \item Organisation is achieved solely through naming conventions.
  \item Related files are grouped using a similar prefix or numbering scheme.
\end{itemize}


\subsection*{Mind Palace Technique}
\textbf{Source:} \cite{metivierMagneticMemoryMethod2024}

Popularised by the fictional Sherlock Holmes, but a real and powerful memorisation technique:

\begin{enumerate}
  \item Imagine a physical location you know well and can vividly visualise.
  \item Assign memories to physical spots within that space.
  \item \textit{Example:} You always hang your keys on the coat hook. Assign the memory of a lock combination to those keys.
  \item Use large, distinct objects (like beds or tables). Avoid duplicating items between rooms.
\end{enumerate}